\begin{afterword}
\summaryhead

The essay proposes that rationality is a martial art of the mind. Just as anyone with a hand can learn to make a fist, anyone with a brain can learn to use it well, and correct its systematic errors. Minds are harder to train than muscles, because reflection is a recent evolutionary development.

The author asks why cities have martial arts schools but no schools of rationality. One reason is that reasoning skill is hard to verify: a board breaks visibly, a thought does not. But recent science, the heuristics-and-biases program, Bayesian probability, evolutionary and social psychology, now gives us the knowledge and vocabulary to build such an art. The author says he learned what he knows of rationality from working on artificial general intelligence. The art must work with brain machinery we cannot rewire. Turning science into practice may be awkward, but introspection, though blurred, is real, so we can apply the science to correct our own thinking.

\responsehead

This is a founding document. It names a discipline, gives it a metaphor, declares that ``humanity may finally be ready'' for it, and gives its author's work on artificial intelligence as the source of his understanding. Read as an argument, it has two gaps, and it states the larger of them itself.

First, the analogy fails where it matters. The essay admits that skill in rationality, unlike a broken board, cannot easily be verified. Without verification there is no dojo in any useful sense: no sparring partner, no teacher's fist to copy, no way to tell progress from the feeling of progress. The essay treats this as one possible reason dojos do not exist, and then proposes the dojo anyway. What survives of the analogy is the costume: the dojo, the levels, the training.

Second, the gap it proposes to fill was not empty. Logic and mathematics have been taught for centuries, and in 1980 the California State University system required formal instruction in critical thinking on all its campuses, for about 300,000 students. The essay does not mention any of this. Its question, ``Why aren't there dojos that teach rationality?'', works only on a reader who does not know about them. The list of new sciences meant to found the art is thin as well: no finding from evolutionary or social psychology is given.

The context sharpens the first point. The Center for Applied Rationality was founded in 2012 out of the LessWrong community that grew from this writing. It now charges \$1,900 to \$4,900, on a sliding scale, for a four-and-a-half-day workshop. The evidence of effect cited on its Wikipedia page is a 50-person trial, reported by \textsc{Vice} in 2016, that measured neuroticism and self-efficacy, not reasoning. I infer that the verification problem this essay raised in 2006 was never solved; it was priced.

In short: a manifesto for a school whose skill, by its own admission, cannot be checked, proposed to fill a gap that did not exist.
\end{afterword}
